%% ============================================================
%%  SECCIÓN 9 — Tratamiento de Datos
%% ============================================================

\section{Tratamiento de Datos y Resultados}
\label{sec:tratamiento}

En esta sección se desarrolla el procesamiento analítico y cuantitativo de los datos experimentales, aplicando las ecuaciones del fundamento teórico y las fórmulas de propagación de incertidumbres deducidas en el Anexo~\ref{anexo:propagacion_errores}.

\subsection{Procesamiento de Tensiones Continuas (DC) y Regresión Lineal}
A partir de las mediciones registradas en escala óptima (Tabla~\ref{tab:resultados_dc}), se calcula la tensión medida en el osciloscopio $V_{\text{osc}}$ y su incertidumbre asociada $\delta V_{\text{osc}}$:

\begin{enumerate}
    \item \textbf{Pila Alcalina ($1{,}5\,\unit{\volt}$ nominal, $S_v = \qty{0,2}{\volt\per\division}$):}
    \begin{equation}
        V_{\text{osc}} = 7{,}2\,\text{div} \times \qty{0,2}{\volt\per\division} = \qty{1,440}{\volt}, \qquad \delta V_{\text{osc}} = 0{,}1\,\text{div} \times \qty{0,2}{\volt\per\division} = \qty{0,020}{\volt}
    \end{equation}
    $$e_r = \frac{|1{,}440 - 1{,}442|}{1{,}442} \times 100\% = \qty{0,14}{\percent}$$

    \item \textbf{Fuente Salida 1 ($2{,}5\,\unit{\volt}$ nominal, $S_v = \qty{0,5}{\volt\per\division}$):}
    \begin{equation}
        V_{\text{osc}} = 4{,}9\,\text{div} \times \qty{0,5}{\volt\per\division} = \qty{2,450}{\volt}, \qquad \delta V_{\text{osc}} = 0{,}1\,\text{div} \times \qty{0,5}{\volt\per\division} = \qty{0,050}{\volt}
    \end{equation}
    $$e_r = \frac{|2{,}450 - 2{,}453|}{2{,}453} \times 100\% = \qty{0,12}{\percent}$$

    \item \textbf{Fuente Salida 2 ($4{,}5\,\unit{\volt}$ nominal, $S_v = \qty{2}{\volt\per\division}$):}
    \begin{equation}
        V_{\text{osc}} = 2{,}6\,\text{div} \times \qty{2}{\volt\per\division} = \qty{5,200}{\volt}, \qquad \delta V_{\text{osc}} = 0{,}1\,\text{div} \times \qty{2}{\volt\per\division} = \qty{0,200}{\volt}
    \end{equation}
    $$e_r = \frac{|5{,}200 - 5{,}220|}{5{,}220} \times 100\% = \qty{0,38}{\percent}$$

    \item \textbf{Fuente Salida 3 ($7{,}5\,\unit{\volt}$ nominal, $S_v = \qty{2}{\volt\per\division}$):}
    \begin{equation}
        V_{\text{osc}} = 4{,}0\,\text{div} \times \qty{2}{\volt\per\division} = \qty{8,000}{\volt}, \qquad \delta V_{\text{osc}} = 0{,}1\,\text{div} \times \qty{2}{\volt\per\division} = \qty{0,200}{\volt}
    \end{equation}
    $$e_r = \frac{|8{,}000 - 8{,}030|}{8{,}030} \times 100\% = \qty{0,37}{\percent}$$

    \item \textbf{Fuente Salida 4 ($10{,}5\,\unit{\volt}$ nominal, $S_v = \qty{2}{\volt\per\division}$):}
    \begin{equation}
        V_{\text{osc}} = 5{,}6\,\text{div} \times \qty{2}{\volt\per\division} = \qty{11,200}{\volt}, \qquad \delta V_{\text{osc}} = 0{,}1\,\text{div} \times \qty{2}{\volt\per\division} = \qty{0,200}{\volt}
    \end{equation}
    $$e_r = \frac{|11,200 - 11,150|}{11,150} \times 100\% = \qty{0,45}{\percent}$$
\end{enumerate}

\subsubsection{Ajuste Lineal por Mínimos Cuadrados}
Para contrastar la relación entre $V_{\text{osc}}$ y $V_{\text{mult}}$, se realiza un ajuste por regresión lineal:
\begin{equation}
    V_{\text{osc}} = m \cdot V_{\text{mult}} + b
\end{equation}
Aplicando el método de mínimos cuadrados mediante el script numérico~\cite{scipy, script_graficas}:
\begin{align}
    m &= 1{,}0020 \pm 0{,}0042 \label{eq:pendiente_dc} \\
    b &= -0{,}007 \pm 0{,}028\,\unit{\volt} \label{eq:intercepto_dc} \\
    R^2 &= 0{,}99995 \label{eq:r2_dc}
\end{align}
La pendiente obtenida $m = 1{,}0020$ coincide con el valor teórico ideal $m_{\text{teo}} = 1{,}0000$ dentro de una discrepancia de apenas \qty{0,20}{\percent}, demostrando una calibración vertical lineal altamente exacta.

\subsection{Tratamiento de la Tensión Alterna del Transformador}
A partir de los datos de la Tabla~\ref{tab:resultados_transformador}:

\subsubsection{Voltaje Pico a Pico y Voltaje Pico}
\begin{align}
    V_{pp} &= 3{,}9\,\text{div} \times \qty{5}{\volt\per\division} = \qty{19,50 \pm 0,50}{\volt} \\
    V_p    &= 4{,}8\,\text{div} \times \qty{2}{\volt\per\division} = \qty{9,60 \pm 0,20}{\volt}
\end{align}
El valor pico teórico derivado del multímetro patrón ($V_{\text{rms}} = \qty{6,830 \pm 0,010}{\volt}$) es:
\begin{equation}
    V_{p,\text{teo}} = \sqrt{2} \cdot V_{\text{rms}} = \sqrt{2} \times 6{,}830 = \qty{9,659 \pm 0,014}{\volt} \implies e_r = \frac{|9{,}60 - 9{,}659|}{9{,}659} \times 100\% = \qty{0,62}{\percent}
\end{equation}

\subsubsection{Voltaje Eficaz Calculado}
\begin{equation}
    V_{\text{rms,calc}} = \frac{V_p}{\sqrt{2}} = \frac{9{,}60}{\sqrt{2}} = \qty{6,788 \pm 0,141}{\volt} \approx \qty{6,79 \pm 0,14}{\volt}
\end{equation}
$$e_r = \frac{|6{,}788 - 6{,}830|}{6{,}830} \times 100\% = \qty{0,59}{\percent}$$

\subsubsection{Período y Frecuencia de la Red}
A una base de tiempo de \qty{2}{\milli\second\per\division}:
\begin{align}
    T &= 8{,}3\,\text{div} \times \qty{2}{\milli\second\per\division} = \qty{16,60 \pm 0,20}{\milli\second} = \qty{0,01660 \pm 0,00020}{\second} \\
    f &= \frac{1}{T} = \frac{1}{0{,}01660} = \qty{60,24}{\hertz}
\end{align}
Propagando la incertidumbre de frecuencia mediante $\delta f = f^2 \delta T$:
\begin{equation}
    \delta f = (60{,}24)^2 \times 0{,}00020 = \qty{0,73}{\hertz} \implies f = \qty{60,24 \pm 0,73}{\hertz}
\end{equation}
$$e_r = \frac{|60{,}24 - 60{,}00|}{60{,}00} \times 100\% = \qty{0,40}{\percent}$$

\subsection{Tratamiento de Señales Periódicas del Generador OWON}
Para una amplitud pico común de $V_p = \qty{4,80 \pm 0,20}{\volt}$ a $f = \qty{250}{\hertz}$:

\begin{enumerate}
    \item \textbf{Onda Senoidal:}
    \begin{equation}
        V_{\text{rms,teo}} = \frac{4{,}80}{\sqrt{2}} = \qty{3,394 \pm 0,141}{\volt}, \quad V_{\text{mult}} = \qty{3,420 \pm 0,010}{\volt} \implies e_r = \qty{0,76}{\percent}
    \end{equation}
    \item \textbf{Onda Triangular:}
    \begin{equation}
        V_{\text{rms,teo}} = \frac{4{,}80}{\sqrt{3}} = \qty{2,771 \pm 0,115}{\volt}, \quad V_{\text{mult}} = \qty{2,790 \pm 0,010}{\volt} \implies e_r = \qty{0,68}{\percent}
    \end{equation}
    \item \textbf{Onda Cuadrada:}
    \begin{equation}
        V_{\text{rms,teo}} = V_p = \qty{4,800 \pm 0,200}{\volt}, \quad V_{\text{mult}} = \qty{4,800 \pm 0,020}{\volt} \implies e_r = \qty{0,00}{\percent}
    \end{equation}
\end{enumerate}

Para el período temporal medido a $S_t = \qty{0,5}{\milli\second\per\division}$ ($N_x = 8{,}0 \pm 0{,}1\,\text{div}$):
\begin{align}
    T &= 8{,}0\,\text{div} \times \qty{0,5}{\milli\second\per\division} = \qty{4,00 \pm 0,05}{\milli\second} = \qty{0,00400 \pm 0,00005}{\second} \\
    f &= \frac{1}{0{,}00400} = \qty{250,0}{\hertz}, \qquad \delta f = (250{,}0)^2 \times 0{,}00005 = \qty{3,1}{\hertz} \implies f = \qty{250,0 \pm 3,1}{\hertz}
\end{align}
$$e_r = \frac{|250{,}0 - 250{,}0|}{250{,}0} \times 100\% = \qty{0,00}{\percent}$$
